\section{Research Agenda and Open Questions}\label{sec:research}

We close the technical part of the paper with a structured research
agenda. The agenda is calibrated against the gaps identified in
\Cref{sec:lit_gaps} and against the methodological openings left by the
case studies of \Cref{sec:case_studies}, and is organised along five
branches: theoretical open questions (\Cref{sec:research_theory}),
methodological research avenues (\Cref{sec:research_method}),
computational and algorithmic challenges
(\Cref{sec:research_compute}), domain-specific opportunities
(\Cref{sec:research_domains}), and industry- and implementation-facing
problems (\Cref{sec:research_industry}). Eleven formal Open Questions,
fourteen Research Avenues, and four Computational Challenges are
distributed across these branches; \Cref{fig:research_mindmap}
summarises the overall structure, and each item is formulated to be
specific enough to underwrite a self-contained research project or
grant proposal.

\subsection{Theoretical open questions}\label{sec:research_theory}

The theoretical core of TCDA is the formal study of the four-tuple
$(\Mc, G, \Tc, \Cc)$ of \Cref{def:tcda} and of the conditions under
which its components interact coherently. The questions below organise
the agenda along the algebraic, categorical, statistical, and
interpretive axes of that interaction.

\begin{question}[Causal persistent modules]\label{q:causal-modules}
Let $(P_t)_{t \ge 0}$ be a one-parameter family of interventional
distributions induced by a continuously varying intervention $\doop(x_t)$
on an SCM $\Mc$. Define a \emph{causal persistent module}
$M^{\Cc}_\bullet$ as the persistent homology of $(P_t)_{t \ge 0}$ with
respect to a $\Cc$-graded filtration, so that birth and death scales
encode the strength and duration of causal effects. Does
$M^{\Cc}_\bullet$ admit a Crawley--Boevey-type decomposition into
interval modules indexed by causal effects, and if so, do the resulting
intervals coincide with classically identified causal estimands?
\end{question}

A positive answer would give an algebraic invariant of an SCM whose
indecomposables are causal effects, in direct analogy with the way the
indecomposables of an ordinary persistent module are barcode intervals.
The question is open at every level: existence, uniqueness, stability,
and computability.

\begin{question}[Functorial adjoints for TCDA]\label{q:adjoints}
Let $\Tc : \mathbf{SCM} \to \mathbf{Pers}$ be a topological functor in
the sense of \Cref{def:tcda}. Does $\Tc$ admit a left or right adjoint
$\Tc^\sharp : \mathbf{Pers} \to \mathbf{SCM}$, and if so, what causal
quantities does the unit
$\eta : \mathrm{id} \Rightarrow \Tc^\sharp \circ \Tc$ or the counit
$\epsilon : \Tc \circ \Tc^\sharp \Rightarrow \mathrm{id}$ compute?
\end{question}

Adjoints, if they exist, would systematically convert TDA-level data
into the SCM that ``best explains'' it, providing a categorical
foundation for inverse problems in causal discovery. The question
sharpens the sheaf-theoretic programme of \citet{mahadevan2021causal,
mahadevan2022universal,mahadevan2025sheaf} and ties it to the
functorial-data-analysis tradition of \citet{spivak2014category,
curry2014sheaves,hajij2023topological}.

\begin{question}[Causal stability theorem]\label{q:causal-stability}
Let $\Mc_1, \Mc_2$ be SCMs in the weak Ibeling--Icard topology on
$\mathfrak{M}$ and let $\Tc = \PHk{k}$. Suppose
$\dB\bigl(\Tc(p^{\Mc_1}), \Tc(p^{\Mc_2})\bigr) < \epsilon$. What is the
tightest upper bound, as a function of $\epsilon$ and $G$, on
\[
\sup_{x,y} \bigl\lVert P_{\Mc_1}(y \mid \doop(x))
                    - P_{\Mc_2}(y \mid \doop(x)) \bigr\rVert_{\mathrm{TV}}\,?
\]
\end{question}

This is the quantitative completion of \Cref{prop:stability-transfer}:
the proposition shows that any Lipschitz causal functional inherits
robustness from bottleneck stability, but it leaves open the explicit
bound for the central object of Pearl's framework---the interventional
distribution itself. A positive answer would bridge the sheaf
programme \citep{mahadevan2025sheaf} and the AIPW programme
\citep{kim2026topological,chernozhukov2018double} by giving the
``other half'' of the dictionary: not just that diagrams perturb under
data perturbation, but that do-effects perturb under diagram
perturbation.

\begin{question}[Identifiability of $G$ from a topological functor]
\label{q:tcda-identifiability}
For a fixed admissible topological functor $\Tc$ and a class
$\mathcal{G}$ of DAGs, characterise necessary and sufficient
conditions for the assignment $G \mapsto \Tc(\mathcal{P}_G)$ to be
injective on $\mathcal{G}$. In particular: for which classes
$\mathcal{G}$ is $\Tc = \PHk{1}$ of the windowed dependence network of
\citet{elyaagoubi2023statistical} identifying?
\end{question}

\Cref{prop:identifiability} reduces TCDA identifiability to an
injectivity question; \Cref{q:tcda-identifiability} asks for the
algebraic characterisation of that injectivity. The classical
identifiability results for linear non-Gaussian models
\citep{shimizu2006lingam} and for score-based topological orderings
\citep{rolland2022score,tran2024constraining,xu2025information,
nocurl2024,wu2026optimization} are the natural touchstones.

\begin{question}[Causal semantics for persistent cycles]
\label{q:cycle-semantics}
Develop a calculus for assigning a causal label
(\emph{feedback}, \emph{confounding}, \emph{mediation}, \emph{spurious})
to a persistent $1$-cycle in $\Dgmk{1}(\Tc(X))$. Specifically, give
sufficient observational and interventional tests that distinguish a
cycle generated by a directed feedback loop in $G$ from a cycle
generated by an unobserved low-dimensional confounder.
\end{question}

This is the formal language whose absence is identified as a gap in
\Cref{sec:lit_gaps}. \Cref{prop:confounding_signature} supplies a
one-sided test (residual cycle signals possible confounding); the
question asks for the analogous tests for the other three cases and
for the cohomological-charge calculus that would tie them together.

\begin{question}[Persistent cohomology for mediation]
\label{q:cohomology-mediation}
Let $A \to M \to Y$ be a mediation triple with a multidimensional
mediator $M \in \R^d$. Apply the Hodge decomposition
\citep{desilva2007coverage,nocurl2024} to a $1$-chain on the simplicial
complex generated by the empirical distribution of $(A, M, Y)$ and
decompose it into gradient (acyclic), curl (cyclic), and harmonic
(globally non-trivial) components. Does the resulting decomposition
recover the Pearl natural direct/natural indirect effect decomposition
\citep{pearl2009causality}, and, if so, under what regularity
conditions on the mediator distribution?
\end{question}

A positive answer would convert mediation analysis from a parametric
exercise to a cohomological computation; the harmonic component, in
particular, would quantify mediation that no local mediator captures.

\begin{question}[Sheaf-cohomological reading of AIPW stability]
\label{q:sheaf-aipw}
Re-derive the stability bound of \citet[Theorem on silhouette
stability]{kim2026topological} as a sheaf-cohomological obstruction in
the topos of \citet{mahadevan2025sheaf}. Identify the obstruction class
that vanishes if and only if the AIPW estimator on $\Tc$ inherits
$\sqrt n$ rates under product-rate nuisance conditions.
\end{question}

\Cref{q:sheaf-aipw} is the most concrete instance of the
methodological dictionary called for in \Cref{sec:lit_gaps}: the
categorical and the asymptotic-statistical halves of TCDA must speak
the same language before the field can be regarded as mature.

\begin{question}[Functional CLT for TCDA beyond silhouettes]
\label{q:functional-clt}
Establish a functional central limit theorem for AIPW-type estimators
of $\Cc = \E[\Tc(Y^1) - \Tc(Y^0)]$ for general admissible $\Tc$,
including persistence landscapes \citep{bubenik2015statistical},
persistence images \citep{adams2017persistence}, and Euler
characteristic curves. Identify the natural function space and the
covariance structure of the limit.
\end{question}

The silhouette case is settled by \citet{kim2026topological}, but
landscapes, images, and Euler curves carry different smoothness and
sparsity properties and require their own central-limit machinery,
parallel to but distinct from \citet{bubenik2015statistical}'s
functional CLT for landscape means.

\begin{question}[Positivity for non-Euclidean outcomes]
\label{q:positivity-noneuclidean}
Reformulate the positivity assumption (C3) of \Cref{sec:prelim_ci} for
outcomes valued in a non-Euclidean topological space $\Mc$. Is the
classical pointwise overlap condition
$0 < \Prob(A = 1 \mid X) < 1$ sufficient, or is a topological
covering condition---e.g.\ each connected component of $\mathrm{supp}\,P^X$
intersects both treatment arms---also required for the topological
contrasts of \cref{eq:tate_silhouette,eq:climate_contrast} to be
identifiable?
\end{question}

Positivity violations occur silently when the outcome lives on a
manifold or stratified space: a covariate stratum can have positive
probability on one connected component of $\Mc$ and zero on another,
without violating the classical scalar overlap.
\Cref{q:positivity-noneuclidean} asks for the topological refinement
of overlap that detects this.

\begin{question}[Transportability across topological domains]
\label{q:transportability}
Develop a transportability theory \citep{pearl2014external,
bareinboim2020pearls} for TCDA estimators: under what conditions on
$\Tc$, on the imaging or measurement modality, and on the causal
graph $G$ can $\widehat\Cc_n$ estimated on a source domain
$(\Mc^{\mathrm{src}}, G)$ be transported to a target domain
$(\Mc^{\mathrm{tgt}}, G)$ with bounded loss?
\end{question}

The case study of \Cref{sec:cs_tate} motivates the question: the
SARS-CoV-2 CT-scan TATE of \citet{kim2026topological} is evaluated on
a single scanner protocol, and the practical deployment of a TCDA
estimator across heterogeneous clinical sites requires a
transportability calculus.

\begin{question}[Topological do-calculus]\label{q:topo-docalculus}
Develop a topological do-calculus: a set of rewrite rules acting on
persistence diagrams, landscapes, or silhouettes that mirror Pearl's
three do-calculus rules and produce, in the limit, the same
identification answers. Identify the topological analogues of the
three Pearlian conditions (insertion/deletion of observations,
exchanging actions and observations, and insertion/deletion of
actions).
\end{question}

A topological do-calculus would extend the categorical work of
\citet{mahadevan2025sheaf} into a constructive algorithmic procedure,
and would supply the inferential engine that
\Cref{q:tcda-identifiability} demands.

\subsection{Methodological research avenues}\label{sec:research_method}

The methodological agenda translates the theoretical questions of
\Cref{sec:research_theory} into algorithms, estimators, and learning
procedures. We enumerate five core avenues.

\begin{research}[Causal Mapper]\label{r:causal-mapper}
Design a \emph{causal} variant of the Mapper algorithm
\citep{singh2007mapper,carriere2018statistical} whose output is not a
simplicial complex but a partially directed mixed graph (PDMG) over
Mapper nodes, with edge orientations inferred from local conditional
independence tests within each cluster. Provide bootstrap-based
confidence guarantees on the orientation of each edge, building on
\citet{lecci2014statistical}.
\end{research}

The case study of \Cref{sec:cs_mapper} sketches the calibrated-CATE
side of this avenue; \Cref{r:causal-mapper} asks for the algorithmic
upgrade that recovers a graph rather than a complex.

\begin{research}[Persistent causal homotopy]\label{r:causal-homotopy}
Develop a theory of \emph{persistent causal homotopy} for time-varying
causal systems: a two-parameter persistence framework
\citep{carlsson2009theory,lesnick2015interactive} in which one
direction is the topological filtration scale and the other is causal
time, and the bigraded module encodes both the topology of cross-
sectional dependence and its temporal evolution.
\end{research}

\citet{elyaagoubi2023statistical} provide the empirical analogue
(windowed dependence networks); \Cref{r:causal-homotopy} promotes the
analogue to a calibrated bigraded estimator with formal interpretation
in terms of regime changes and causal-structure shifts.

\begin{research}[Topological regularisation for neural causal models]
\label{r:topo-regularisation}
Use topological summaries as regularisers in deep generative causal
models, extending \citet{xu2021topological,farzam2025topology,
guan2025causal,ma2025causal}. The objective is a causal variational
autoencoder whose latent space is constrained to have prescribed
persistent homology, encouraging causally meaningful disentanglement.
Provide identifiability guarantees and downstream causal-effect
estimation rates.
\end{research}

\begin{research}[Multi-scale causal inference]\label{r:multiscale}
Treat the filtration parameter $t$ of $\PHk{k}(\Filt)$ as a
\emph{causal scale}: at fine scales the topology resolves
micro-mechanisms; at coarse scales it resolves macro-mechanisms.
Develop estimators of $\Cc$ as functions of $t$, with stability and
identifiability guarantees that interpolate between the two regimes.
\end{research}

This avenue formalises an intuition long present in TDA---multi-scale
structure is the field's defining property---in a causal language. The
case study of \Cref{sec:cs_mapper} motivates the avenue: regime-
conditional CATEs naturally depend on a scale.

\begin{research}[Topological mediation via Hodge decomposition]
\label{r:hodge-mediation}
Implement the mediation decomposition of \Cref{q:cohomology-mediation}
as an estimator: compute the empirical Hodge decomposition of a $1$-chain
on the data simplicial complex, identify the gradient/curl/harmonic
components, and map each to a causal estimand (direct, indirect, and
non-local effect) with bootstrap confidence sets.
\end{research}

\subsection{Computational and algorithmic challenges}\label{sec:research_compute}

Beyond statistical and mathematical theory, four computational
problems gate the practical deployment of TCDA at the scale of modern
scientific data.

\begin{challenge}[Scalable persistent homology for causal-graph contexts]
\label{c:scalable-ph}
Existing persistent-homology libraries \citep{bauer2021ripser,
gudhi2021,tauzin2021giotto} compute $\PHk{k}$ for $k \le 2$ in
$O(n^3)$ practical time for $n \lesssim 10^5$ points. Causal-graph
applications routinely require $n \gtrsim 10^6$ (population-level
omics, internet-scale finance). Develop sparse, approximate, or
distributed PH variants whose accuracy guarantees suffice for the
stability bound of \Cref{prop:stability-transfer} to remain
informative at the resulting noise level.
\end{challenge}

\begin{challenge}[Probabilistic persistent homology]\label{c:bayesian-ph}
Calibrated uncertainty in TCDA requires probabilistic versions of
$\PHk{k}$ \citep{maroulas2020bayesian,fasy2014confidence}. Develop
Bayesian inference over persistence diagrams that scales beyond the
small-$n$ regime currently feasible and whose posterior contracts at
the rate established by \citet{lecci2014statistical} for the bootstrap.
\end{challenge}

\begin{challenge}[Interoperable software for TCDA]\label{c:software}
The TDA stack (\texttt{Ripser}, \texttt{GUDHI}, \texttt{giotto-tda})
and the causal-inference stack (\texttt{DoWhy}, \texttt{causal-learn},
\texttt{EconML}, \texttt{CausalML}) are mutually opaque: there is no
common data structure for a persistence diagram with attached causal-
graph metadata, and no shared API for AIPW on a topological summary.
Build a \texttt{tcda-py} (or analogous) library that exposes the TCDA
tuple $(\Mc, G, \Tc, \Cc)$ as a first-class object and orchestrates
estimation, inference, and reporting end-to-end.
\end{challenge}

\begin{challenge}[GPU-accelerated cubical persistent homology]
\label{c:gpu-cubical}
Volumetric medical imaging (3D MR, 4D CT) and gridded climate fields
require cubical $\PHk{k}$ on grids of size $\gtrsim 10^9$ voxels.
Current cubical-PH implementations are CPU-bound and intractable at
this scale. Develop GPU-accelerated, distributed cubical PH whose
output is bit-for-bit compatible with the AIPW pipelines of
\citet{kim2026topological}.
\end{challenge}

\subsection{Domain-specific research opportunities}\label{sec:research_domains}

We collect five domain-specific avenues. Each is anchored in a case
study of \Cref{sec:case_studies} and stated as a substantive,
fundable programme of work.

\paragraph{Healthcare and drug discovery.}
\begin{research}[Topological causal models for drug repurposing]
\label{r:healthcare-repurposing}
Use persistent homology of disease-pathway networks and protein-
ligand binding pockets \citep{kovacev2016using,cang2018integration,
axelrod2022geom} as the topological functor in a TCDA pipeline whose
causal query is the predicted treatment effect of an off-label drug.
The pipeline must produce a calibrated point estimate of efficacy
together with a topological diagnostic for unmeasured comorbidities,
generalising the \citet{kim2026topological} GEOM-Drugs case study.
\end{research}

\paragraph{Climate and Earth-system science.}
\begin{research}[Topological attribution of extreme climate events]
\label{r:climate-attribution}
Build a fully-specified TCDA pipeline that ingests an event-based
climate ensemble (factual + counterfactual scenarios), computes
cubical $\PHk{0,1}$ of the spatial event field, and outputs the
counterfactual contrast of \cref{eq:climate_contrast} with bootstrap
confidence sets. Demonstrate on a benchmark of canonical extreme
events (Western European heatwave 2003, Pakistan flood 2022).
\end{research}

\paragraph{Finance and macroeconomics.}
\begin{research}[Topological causal early-warning systems]
\label{r:finance-warning}
Extend the topological landscape early-warning signals of
\citet{gidea2018topological,ismail2022early,leykam2023topological}
into a causal early-warning framework: detect not only that a regime
shift has occurred but identify the triggering causal factor via a
topological contrast on the underlying network of asset returns,
following \citet{krasnovsky2025causal}.
\end{research}

\paragraph{Neuroscience.}
\begin{research}[Whole-brain causal mapping via topological networks]
\label{r:neuroscience-mapping}
Synthesise the brain-network programme \citep{bullmore2009complex,
sizemore2018cliques,sizemore2019importance,fathian2022trend} with
the simulation-based inference of \citet{elyaagoubi2023statistical}
to deliver whole-brain causal maps in which (i) cycles in the
empirical $\PHk{1}$ are formally tested against null Granger graphs
and (ii) each surviving cycle is annotated with a putative causal
mechanism by the calculus of \Cref{q:cycle-semantics}.
\end{research}

\paragraph{Materials science.}
\begin{research}[Topological causal structure--property relations]
\label{r:materials-structure-property}
Apply the TCDA tuple to materials discovery: $\Mc$ is a configuration
space of atomic positions, $G$ encodes the causal mechanism linking
structure to property, $\Tc = \PHk{0,1,2}$ of the
Vietoris--Rips filtration of atomic positions
\citep{nigmetov2022topological,leykam2023topological}, and $\Cc$ is the
predicted change in property under a topological intervention (e.g.\
introducing a defect). The output is a causal explanation of why a
given microstructure produces a given macroscopic property.
\end{research}

\subsection{Industry and practical implementation}\label{sec:research_industry}

The last four items in our agenda concern the translation of TCDA
research into industrial practice. We state each as a Research Avenue
because they require organised research programmes, not merely
engineering.

\begin{research}[Validation standards for topological causal models]
\label{r:validation-standards}
Develop reporting and validation standards for topological causal
estimators in regulated industries (FDA-relevant clinical decision
support, prudential regulators for financial early-warning systems).
The standards must specify the minimum reportable uncertainty (the
\citet{lecci2014statistical} bootstrap is a candidate baseline), the
stability bound to be reported (a corollary of
\Cref{prop:stability-transfer}), and the confounding diagnostic
threshold (a corollary of \Cref{prop:confounding_signature}).
\end{research}

\begin{research}[Stakeholder interpretability for topological causal findings]
\label{r:interpretability}
Develop interpretability tools that translate topological causal
findings---a persistence diagram, a Mapper graph, a counterfactual
silhouette contrast---into the language of clinicians, policymakers,
and managers. The constructive analogue of the Pearl
``probabilities of causation'' framework
\citep{pearl2009causality,balke1997probabilistic} is a natural
starting point: how does ``a persistent $1$-cycle implies a 75\%
probability of mechanism $M$'' get reported?
\end{research}

\begin{research}[Cloud-native infrastructure for large-scale TCDA]
\label{r:infrastructure}
Build a cloud-native compute layer that exposes the TCDA tuple
$(\Mc, G, \Tc, \Cc)$ as a managed service, with built-in scaling for
the GPU cubical PH of \Cref{c:gpu-cubical}, distributed Vietoris--Rips
of \Cref{c:scalable-ph}, and interoperable wrappers for the libraries
named in \Cref{c:software}. The deliverable is a reference
implementation (open-source) and a managed-service blueprint
(commercial).
\end{research}

\begin{research}[Benchmark datasets for TCDA]\label{r:benchmarks}
Design and release a TCDA benchmark suite analogous to the
CMU-CausalDisco and CausalBench collections
\citep{causaldisco2023survey,timeseriescd2023} but specifically
calibrated to topological-causal interactions: graphs with known
homological signatures, treatment effects with known topological
contrast, confounders with prescribed homotopy type. The simulation
framework of \citet{elyaagoubi2023statistical} is the natural
starting point; the suite must extend to imaging, molecular, and
climate modalities.
\end{research}

% =====================================================================
% Figure: TCDA research-agenda mind map (Section 6).
% Working draft -- the user will refine the final visual styling.
% =====================================================================
\begin{figure}[t]
\centering
\begin{tikzpicture}[
    font=\footnotesize,
    every node/.style={align=center},
    centroid/.style={draw, rounded corners, fill=violet!12, minimum width=42mm,
                     minimum height=10mm, font=\small\bfseries},
    branch/.style={draw, rounded corners, fill=blue!10, minimum width=32mm,
                   minimum height=8mm, font=\footnotesize\bfseries},
    leaf/.style={draw, rounded corners, fill=teal!6, minimum width=34mm,
                 minimum height=6mm, inner sep=2pt, font=\scriptsize},
    edge from parent/.style={draw, semithick, -{Latex[length=2mm]}, gray!75}]

\node[centroid] (root) at (0, 0) {TCDA Research Agenda\\(\Cref{sec:research})};

\node[branch] (theory)   at ($(root) + (-7.4,  3.0)$) {Theoretical\\(\S\ref{sec:research_theory})};
\node[branch] (methods)  at ($(root) + ( 7.4,  3.0)$) {Methodological\\(\S\ref{sec:research_method})};
\node[branch] (compute)  at ($(root) + (-7.4, -3.0)$) {Computational\\(\S\ref{sec:research_compute})};
\node[branch] (domains)  at ($(root) + ( 7.4, -3.0)$) {Domain-specific\\(\S\ref{sec:research_domains})};
\node[branch] (industry) at ($(root) + ( 0.0, -5.4)$) {Industry\\(\S\ref{sec:research_industry})};

\draw[semithick, -{Latex[length=2mm]}, gray!75] (root) -- (theory);
\draw[semithick, -{Latex[length=2mm]}, gray!75] (root) -- (methods);
\draw[semithick, -{Latex[length=2mm]}, gray!75] (root) -- (compute);
\draw[semithick, -{Latex[length=2mm]}, gray!75] (root) -- (domains);
\draw[semithick, -{Latex[length=2mm]}, gray!75] (root) -- (industry);

% Theory leaves (upper-left)
\node[leaf] (t1) at ($(theory) + (-3.0, 1.6)$) {Causal persistent\\modules (Q1)};
\node[leaf] (t2) at ($(theory) + ( 3.0, 1.6)$) {Functorial adjoints\\(Q2)};
\node[leaf] (t3) at ($(theory) + (-3.0, 0.5)$) {Causal stability\\(Q3)};
\node[leaf] (t4) at ($(theory) + ( 3.0, 0.5)$) {Identifiability\\from $\Tc$ (Q4)};
\node[leaf] (t5) at ($(theory) + (-3.0,-0.6)$) {Cycle semantics\\(Q5--Q6)};
\node[leaf] (t6) at ($(theory) + ( 3.0,-0.6)$) {Sheaf--AIPW\\bridge (Q7)};
\foreach \x in {t1,t2,t3,t4,t5,t6} \draw[gray!60] (theory.south) -- (\x.north);

% Methods leaves (upper-right)
\node[leaf] (m1) at ($(methods) + (-3.0, 1.6)$) {Causal Mapper\\(R1)};
\node[leaf] (m2) at ($(methods) + ( 3.0, 1.6)$) {Persistent\\causal homotopy (R2)};
\node[leaf] (m3) at ($(methods) + (-3.0, 0.5)$) {Topological\\regularisation (R3)};
\node[leaf] (m4) at ($(methods) + ( 3.0, 0.5)$) {Multi-scale\\causal inference (R4)};
\node[leaf] (m5) at ($(methods) + (-3.0,-0.6)$) {Hodge mediation\\(R5)};
\foreach \x in {m1,m2,m3,m4,m5} \draw[gray!60] (methods.south) -- (\x.north);

% Compute leaves (lower-left)
\node[leaf] (c1) at ($(compute) + (-3.0,-1.6)$) {Scalable PH\\for causal graphs (C1)};
\node[leaf] (c2) at ($(compute) + ( 3.0,-1.6)$) {Bayesian PH\\(C2)};
\node[leaf] (c3) at ($(compute) + (-3.0,-0.6)$) {Interoperable\\software (C3)};
\node[leaf] (c4) at ($(compute) + ( 3.0,-0.6)$) {GPU cubical PH\\(C4)};
\foreach \x in {c1,c2,c3,c4} \draw[gray!60] (compute.north) -- (\x.south);

% Domain leaves (lower-right)
\node[leaf] (d1) at ($(domains) + (-3.0,-0.6)$) {Healthcare\\(R6)};
\node[leaf] (d2) at ($(domains) + ( 3.0,-0.6)$) {Climate\\(R7)};
\node[leaf] (d3) at ($(domains) + (-3.0,-1.6)$) {Finance\\(R8)};
\node[leaf] (d4) at ($(domains) + ( 3.0,-1.6)$) {Neuroscience\\(R9)};
\node[leaf] (d5) at ($(domains) + ( 0.0,-2.6)$) {Materials\\(R10)};
\foreach \x in {d1,d2,d3,d4,d5} \draw[gray!60] (domains.north) -- (\x.south);

% Industry leaves (bottom)
\node[leaf] (i1) at ($(industry) + (-5.0,-1.2)$) {Validation\\standards (R11)};
\node[leaf] (i2) at ($(industry) + (-1.8,-1.2)$) {Stakeholder\\interpretability (R12)};
\node[leaf] (i3) at ($(industry) + ( 1.8,-1.2)$) {Cloud-native\\infrastructure (R13)};
\node[leaf] (i4) at ($(industry) + ( 5.0,-1.2)$) {Benchmark\\datasets (R14)};
\foreach \x in {i1,i2,i3,i4} \draw[gray!60] (industry.north) -- (\x.south);

\end{tikzpicture}
\caption{Mind-map of the TCDA research agenda. The agenda has five
branches---theoretical (\Cref{sec:research_theory}), methodological
(\Cref{sec:research_method}), computational
(\Cref{sec:research_compute}), domain-specific
(\Cref{sec:research_domains}), and industry-facing
(\Cref{sec:research_industry})---comprising eleven formal Open
Questions (Q1--Q11), fourteen Research Avenues (R1--R14), and four
Computational Challenges (C1--C4). Each leaf is keyed to the formal
statement in the body of \Cref{sec:research}; together the leaves
populate every domain of the TCDA taxonomy (\Cref{fig:taxonomy}) with
at least one concrete, grant-scale open problem.}
\label{fig:research_mindmap}
\end{figure}


\subsection{Synthesis: a five-year and a ten-year horizon}
\label{sec:research_synth}

\Cref{fig:research_mindmap} aggregates the agenda. On a five-year
horizon, the most actionable items are the \emph{computational}
challenges of \Cref{sec:research_compute}---a working
\texttt{tcda-py} library, GPU cubical PH, scalable Vietoris--Rips---
together with the \emph{methodological} avenues of
\Cref{sec:research_method} that are within reach of existing
estimation theory: causal Mapper (\Cref{r:causal-mapper}),
topological regularisation (\Cref{r:topo-regularisation}), and Hodge
mediation (\Cref{r:hodge-mediation}). On a ten-year horizon, the
field's foundational identity will be settled by the \emph{theoretical}
questions of \Cref{sec:research_theory}: causal persistent modules
(\Cref{q:causal-modules}), functorial adjoints
(\Cref{q:adjoints}), a quantitative causal stability theorem
(\Cref{q:causal-stability}), the identifiability characterisation
(\Cref{q:tcda-identifiability}), and the sheaf-cohomological reading
of stability (\Cref{q:sheaf-aipw}). The domain-specific opportunities
of \Cref{sec:research_domains} and the industry items of
\Cref{sec:research_industry} are the empirical loci where the field
will earn its scientific standing.

Two observations close the section. First, every open question and
research avenue connects back to a gap identified in
\Cref{sec:lit_gaps} or to a remaining limitation noted in the case
studies of \Cref{sec:case_studies}: the agenda is not invented from
first principles but is the systematic completion of the existing
literature surveyed in \Cref{sec:literature}. Second, the agenda is
deliberately inclusive in scope and parsimonious in commitments: it
declares which problems matter without prescribing which solutions
will succeed, and so leaves room for a community of researchers from
both TDA and causal inference to enter the field on their own terms.
The next section---\Cref{sec:conclusion}---returns to the overall
vision and to the call for interdisciplinary collaboration that the
agenda enables.
